% Section 4. Drafted from docs/PAPER-OUTLINE.md section 4 and docs/SUBSUMPTION.md. Every named
% result is a gated, axiom-free theorem in coq/*.v; statements are read from the source.

\section{CRDTs, invariant confluence, and what repair adds}\label{sec:crdt}

Two bodies of work answer neighboring questions. Conflict-free replicated data types (CRDTs)
obtain convergence by designing operations so that order never matters~\cite{shapiro2011}.
Invariant confluence asks which transactions can run without coordination while preserving an
invariant~\cite{bailis2014}. This section places normalization confluence against both, precisely:
CRDTs are the case where repair does nothing (\cref{sec:crdt-identity}); repair enlarges which
operations can be used, but not which behaviors converge (\cref{sec:crdt-readings}); and invariant
confluence is a different question, because there the invariant is checked and never repaired
(\cref{sec:crdt-iconf}).

% ---------------------------------------------------------------------------------------------
\subsection{CRDTs are the case where repair does nothing}\label{sec:crdt-identity}

In the terms of \cref{sec:main}, a CRDT is a system whose repair $N$ is the identity. State descent,
$N(e(N(x))) = N(e(x))$, then holds for every event, so convergence comes down to history descent
alone: every reordering the delivery model allows must give the same state. The classical CRDT
conditions are what this reduces to under the two delivery models CRDTs are built for. Write
$\mathrm{gstep}(e) = N \circ e$ for the governed step.

\begin{theorem}[The compensation-free boundary; \coq{crdt_boundary} (a), (b), (b')]\label{thm:boundary}
Let $N$ be the identity.
\begin{enumerate}[label=(\alph*)]
  \item \emph{Causal exactly-once delivery.} For every irreflexive happens-before relation, the
    governed steps converge causally from every start (any two causally consistent orders of the
    same events reach the same state) if and only if concurrent raw operations commute at every
    state.
  \item \emph{Unordered at-least-once delivery.} For every event range $R$ with decidable equality
    and every start $s_0$, any two deliveries over $R$ with the same set of events, in any order and
    with any duplication, reach the same state from $s_0$ if and only if the raw operations are
    commutative and idempotent at every state reachable from $s_0$.
  \item[(b$'$)] \emph{The semilattice form.} If moreover $R$ is finite and state equality is
    decidable, the convergence of (b) holds if and only if the raw operations have a
    join-semilattice representation on the states reachable from $s_0$: some join is a semilattice
    there and every event acts as $e(s) = s \sqcup e(s_0)$.
\end{enumerate}
\end{theorem}

Part (a) is the op-based (CmRDT) condition and part (b$'$) the state-based (CvRDT) condition of
Shapiro et al.~\cite{shapiro2011}, each in exact form: Shapiro et al.\ give them as sufficient
conditions, and the theorem shows they are also necessary, on the states a system can reach. The
qualifiers are needed: reachability cannot be dropped (\coq{clamp_reach_qualifier}), a semilattice on
all of the state space is too strong (\coq{naive_cvrdt_iff_fails}), and both commutativity and
idempotence are needed in (b) (\coq{lww_not_conv}, \coq{add_cc_not_alo}). The exact conditions behind
(a) and (b) are \coq{causal_convergence_exact} and \coq{merge_action_exact}; the theorem instantiates
them at the identity repair.

% ---------------------------------------------------------------------------------------------
\subsection{Two readings of ``repair enlarges what converges''}\label{sec:crdt-readings}

The usual claim for a repair-based design is that it reaches systems CRDTs cannot. The claim is true
in one reading and false in another, and the paper separates them.

\paragraph{Operations: strict.}
Some raw operations converge only with repair.

\begin{proposition}[\coq{crdt_boundary} (c)]\label{prop:witness}
Take a single Boolean whose only valid value is \textsf{false}, two events with raw transitions
$\mathsf{settrue}(s) = \textsf{true}$ and $\mathsf{flip}(s) = \neg s$, and the repair that resets to
\textsf{false}. The governed system converges under causal exactly-once delivery from every start
and under unordered at-least-once delivery from every start. The raw operations satisfy none of the
conditions of \cref{thm:boundary}: they do not commute where the two events are concurrent, the
system built from them without repair fails to converge under both delivery models, they are neither
commutative nor idempotent at a reachable state from any start, and they have no semilattice
representation on the reachable states.
\end{proposition}

So no system without repair that uses these raw transitions converges, while the repaired system
does. The separation is sharpened in \coq{witness_ops_not_commute}, \coq{witness_raw_not_causal},
\coq{witness_not_cvrdt_order} (no partial order on the Booleans makes both raw operations
inflationary) and \coq{witness_not_cvrdt_exact}.

\paragraph{Behavior: not strict.}
The exact conditions behind \cref{thm:boundary} hold for any step function, not only for raw
operations. Read at the governed step $\mathrm{gstep}(e) = N \circ e$, they say:

\begin{corollary}[\coq{causal_convergence_exact}, \coq{merge_action_exact}, at the governed step]\label{cor:behavior}
For any repair $N$:
\begin{enumerate}[label=(\alph*)]
  \item for an irreflexive happens-before relation, the governed system converges causally from
    every start if and only if governed steps of concurrent events commute at every state;
  \item from a start $s_0$, the governed system converges under unordered at-least-once delivery if
    and only if the governed steps are commutative and idempotent at every state reachable from
    $s_0$.
\end{enumerate}
\end{corollary}

These are the CmRDT and the commutative-idempotent conditions, stated for the governed steps. So
every convergent governed system is a CRDT on its own governed steps, at the states it reaches. The
witness of \cref{prop:witness} is an instance: its governed behavior is the constant
\textsf{false}, and its governed steps are a CvRDT with join $\wedge$
(\coq{witness_governed_constant}).

\paragraph{What repair adds.}
Normalization confluence does not reach convergent behaviors that CRDTs cannot. What it adds is a
way to build them. A designer writes the operations the application needs, which need not commute,
and a repair that restores the invariants; the exact conditions of \cref{sec:main}, state descent
and history descent, decide whether the combination converges, which is to say whether its governed
steps form a CRDT. The design freedom is in the operations, and the guarantee is checked rather than
designed in. \cref{sec:gsm} describes a checker that runs these conditions at build time.

% ---------------------------------------------------------------------------------------------
\subsection{Invariant confluence: checked, not repaired}\label{sec:crdt-iconf}

Invariant confluence~\cite{bailis2014} asks whether a set of transactions can run on replicas that
merge their states without coordination and still preserve an invariant $I$. A transaction commits
only if its local result satisfies $I$, and a merge can never be refused. The invariant is checked,
never repaired. This is a different question from the one normalization confluence answers, and the
two meet in one place: the merge.

Model the invariant as the solution set of a constraint problem and the merge as an operation
applied item by item. In the constraint-satisfaction terms of \cref{sec:layer-b}, the merge is then
a binary operation on states, and invariant confluence becomes a statement about it.

\begin{theorem}[\coq{pol_iconfluent}, \coq{iconf_iff_pol_reach}]\label{thm:iconf}
\leavevmode
\begin{enumerate}[label=(\alph*)]
  \item If the item-wise merge is a polymorphism of the invariant's constraint language, every set of
    transactions is invariant confluent.
  \item Exactly: when $I$ holds at the start and reads finitely many items, a set of transactions is
    invariant confluent if and only if the merge is a polymorphism of the relation of states
    reachable under those transactions with $I$ preserved.
\end{enumerate}
\end{theorem}

The restriction to reachable states is necessary, as Whittaker and
Hellerstein~\cite{whittaker2018} observed. A mechanized counterexample makes it concrete
(\coq{uq_not_pol}): with two versions $a$ and $b$, merge by set union, the invariant ``not both
present'' and the single transaction ``add $a$'', the transactions are invariant confluent, but
union does not preserve the invariant on all states. Bailis et al.'s main theorem, that a
coordination-free, available, convergent and invariant-preserving execution exists exactly for
invariant-confluent transactions, holds in a replica model mechanized in the same file
(\coq{replica_exact}); necessity for arbitrary protocols would need a model of partitions that the
development does not have.

Repair changes the picture. With the same invariant and both transactions ``add $a$'' and ``add
$b$'', the transactions are not invariant confluent. With the repair
$N(x, y) = (x,\, y \wedge \neg x)$, which drops $b$ whenever $a$ is present, state descent and
history descent both hold, every repaired state satisfies the invariant, and replicas that exchange
repaired states agree with the repaired run of the joint history (\coq{repair_enlarges}). The caveat
is that this repair can undo a committed write, which Bailis et al.'s model forbids; the comparison is
between two models, not a strengthening of theirs.
